% GENERATED FILE. Edit canonical lessons, then run npm run generate:books.
% canonical-source-hash: fnv1a64:abda5271600f13bb
% canonical-lessons: SA-C134-chinatti, SA-C134-bhinatti, SA-C134-sthapayati, SA-C134-presayati, SA-C134-anvisyati

\chapter{Cuts, Breaks, Places, Sends, Searches}
\label{ch:sa-cuts-breaks-places-sends-searches}
\hlchaptermodality{134}

\begin{chapteropening}
\textbf{By the end of this chapter:} \emph{I can say cuts, breaks, places, sends and searches.}

Complete the last lesson of chapter 134: i can say cuts, breaks, places, sends and searches.
\end{chapteropening}

\section[chinatti]{\sk{छिनत्ति} (chinatti) — cuts}

\label{lesson:SA-C134-chinatti}



\begin{quote}
Before the new one: say the Sanskrit for swims, then the Sanskrit for jumps.
\end{quote}

\subsection*{You'll want to know: \sk{छिनत्ति}}
\textbf{\sk{छिनत्ति}} — \emph{chinatti} — \textquotedblleft{}cuts\textquotedblright{}.

\begin{grammarlens}[title={What you've built}]
One more word for this chapter.
\end{grammarlens}

\subsection*{Your turn}
\begin{itemize}
\raggedright
  \item \emph{Say it:} \emph{chinatti}
  \item \emph{Say it:} \emph{chinatti}, once more
  \item \emph{Say it:} say \emph{chinatti}
  \item \emph{Recall:} say the Sanskrit for swims, then the Sanskrit for jumps, then say \emph{chinatti} again
\end{itemize}

\begin{tcolorbox}[breakable,colback=teal!4,colframe=teal!35!black,title={Before you move on}]
What does \sk{छिनत्ति} mean? (\textquotedblleft{}Cuts\textquotedblright{}.) Say it once more.
\end{tcolorbox}
\section[bhinatti]{\sk{भिनत्ति} (bhinatti) — breaks}

\label{lesson:SA-C134-bhinatti}



\begin{quote}
Before the new one: say the Sanskrit for jumps, then the Sanskrit for cuts.
\end{quote}

\subsection*{You'll want to know: \sk{भिनत्ति}}
\textbf{\sk{भिनत्ति}} — \emph{bhinatti} — \textquotedblleft{}breaks\textquotedblright{}.

\begin{grammarlens}[title={What you've built}]
One more word for this chapter.
\end{grammarlens}

\subsection*{Your turn}
\begin{itemize}
\raggedright
  \item \emph{Say it:} \emph{bhinatti}
  \item \emph{Say it:} \emph{bhinatti}, once more
  \item \emph{Say it:} say \emph{bhinatti}
  \item \emph{Recall:} say the Sanskrit for jumps, then the Sanskrit for cuts, then say \emph{bhinatti} again
\end{itemize}

\begin{tcolorbox}[breakable,colback=teal!4,colframe=teal!35!black,title={Before you move on}]
What does \sk{भिनत्ति} mean? (\textquotedblleft{}Breaks\textquotedblright{}.) Say it once more.
\end{tcolorbox}
\section[sthāpayati]{\sk{स्थापयति} (sthāpayati) — places, puts}

\label{lesson:SA-C134-sthapayati}



\begin{quote}
Before the new one: say the Sanskrit for cuts, then the Sanskrit for breaks.
\end{quote}

\subsection*{You'll want to know: \sk{स्थापयति}}
\textbf{\sk{स्थापयति}} — \emph{sthāpayati} — \textquotedblleft{}places, puts\textquotedblright{}.

\begin{grammarlens}[title={What you've built}]
One more word for this chapter.
\end{grammarlens}

\subsection*{Your turn}
\begin{itemize}
\raggedright
  \item \emph{Say it:} \emph{sthāpayati}
  \item \emph{Say it:} \emph{sthāpayati}, once more
  \item \emph{Say it:} say \emph{sthāpayati}
  \item \emph{Recall:} say the Sanskrit for cuts, then the Sanskrit for breaks, then say \emph{sthāpayati} again
\end{itemize}

\begin{tcolorbox}[breakable,colback=teal!4,colframe=teal!35!black,title={Before you move on}]
What does \sk{स्थापयति} mean? (\textquotedblleft{}Places, puts\textquotedblright{}.) Say it once more.
\end{tcolorbox}
\section[preṣayati]{\sk{प्रेषयति} (preṣayati) — sends}

\label{lesson:SA-C134-presayati}



\begin{quote}
Before the new one: say the Sanskrit for breaks, then the Sanskrit for places.
\end{quote}

\subsection*{You'll want to know: \sk{प्रेषयति}}
\textbf{\sk{प्रेषयति}} — \emph{preṣayati} — \textquotedblleft{}sends\textquotedblright{}.

\begin{grammarlens}[title={What you've built}]
One more word for this chapter.
\end{grammarlens}

\subsection*{Your turn}
\begin{itemize}
\raggedright
  \item \emph{Say it:} \emph{preṣayati}
  \item \emph{Say it:} \emph{preṣayati}, once more
  \item \emph{Say it:} say \emph{preṣayati}
  \item \emph{Recall:} say the Sanskrit for breaks, then the Sanskrit for places, then say \emph{preṣayati} again
\end{itemize}

\begin{tcolorbox}[breakable,colback=teal!4,colframe=teal!35!black,title={Before you move on}]
What does \sk{प्रेषयति} mean? (\textquotedblleft{}Sends\textquotedblright{}.) Say it once more.
\end{tcolorbox}
\section[anviṣyati]{\sk{अन्विष्यति} (anviṣyati) — searches}

\label{lesson:SA-C134-anvisyati}



\begin{quote}
Before the new one: say the Sanskrit for places, then the Sanskrit for sends.
\end{quote}

\subsection*{You'll want to know: \sk{अन्विष्यति}}
\textbf{\sk{अन्विष्यति}} — \emph{anviṣyati} — \textquotedblleft{}searches\textquotedblright{}.

\begin{grammarlens}[title={What you've built}]
One more word for this chapter.
\end{grammarlens}

\subsection*{Your turn}
\begin{itemize}
\raggedright
  \item \emph{Say it:} \emph{anviṣyati}
  \item \emph{Say it:} \emph{anviṣyati}, once more
  \item \emph{Say it:} say \emph{anviṣyati}
  \item \emph{Recall:} say the Sanskrit for places, then the Sanskrit for sends, then say \emph{anviṣyati} again
\end{itemize}

\begin{tcolorbox}[breakable,colback=teal!4,colframe=teal!35!black,title={Before you move on}]
What does \sk{अन्विष्यति} mean? (\textquotedblleft{}Searches\textquotedblright{}.) Say it once more.
\end{tcolorbox}
